\documentclass[../../main.tex]{subfiles}
\begin{document}

{\bf [解]}
各点の位置ベクトルを対応する小文字で表す．題意より
\begin{align}
\vec k=\dfrac{\vec a+\vec b}{2},\quad \vec l=\dfrac{\vec b+\vec c}{2},\quad \vec m=\dfrac{\vec c+\vec d}{2},\quad \vec n=\dfrac{\vec d+\vec a}{2} \label{eq:1}
\end{align}
である．

\begin{figure}[htb]
\centering
\begin{tikzpicture}[scale=1.2]
  % 四面体ABCDの模式的な配置 (実際の3次元形状を厳密には反映しない見取り図)
  \coordinate (A) at (0,2.2);
  \coordinate (B) at (-1.6,-0.3);
  \coordinate (C) at (1.0,-0.6);
  \coordinate (D) at (2.0,0.6);
  \draw (A)--(B)--(C)--cycle;
  \draw (A)--(D);
  \draw[dashed] (B)--(D);
  \draw (C)--(D);
  % 各辺の中点 K,L,M,N
  \coordinate (K) at ($(A)!0.5!(B)$);
  \coordinate (L) at ($(B)!0.5!(C)$);
  \coordinate (M) at ($(C)!0.5!(D)$);
  \coordinate (N) at ($(D)!0.5!(A)$);
  \draw[dotted] (K)--(L)--(M)--(N)--cycle;
  \filldraw (A) circle(1.0pt) node[above]{$A$};
  \filldraw (B) circle(1.0pt) node[left]{$B$};
  \filldraw (C) circle(1.0pt) node[below]{$C$};
  \filldraw (D) circle(1.0pt) node[right]{$D$};
  \filldraw (K) circle(1.0pt) node[left]{$K$};
  \filldraw (L) circle(1.0pt) node[below]{$L$};
  \filldraw (M) circle(1.0pt) node[right]{$M$};
  \filldraw (N) circle(1.0pt) node[above right]{$N$};
\end{tikzpicture}
\caption{四面体$ABCD$と各辺の中点$K,L,M,N$}
\end{figure}

\paragraph{(1)}

\cref{eq:1}を用いて題意の内積を計算すると

\begin{align}
\overrightarrow{MK}\cdot\overrightarrow{LN}
  &=(\vec k-\vec m)\cdot(\vec n-\vec l) \\
  &=\dfrac{1}{4}(\vec a+\vec b-\vec c-\vec d)\cdot(\vec a+\vec d-\vec b-\vec c)\\
  &=\dfrac{1}{4}\bigl[(\vec a-\vec c)+(\vec b-\vec d)\bigr]\cdot\bigl[(\vec a-\vec c)-(\vec b-\vec d)\bigr]\\
  &=\dfrac{1}{4}\bigl(|\vec a-\vec c|^2-|\vec b-\vec d|^2\bigr) \\
  &=\dfrac{1}{4}\bigl(|\overrightarrow{AC}|^2-|\overrightarrow{BD}|^2\bigr) \\
\therefore 
  4\overrightarrow{MK}\cdot\overrightarrow{LN}
  &\bigl(|\overrightarrow{AC}|^2-|\overrightarrow{BD}|^2\bigr)
\end{align}
となり題意は示された．

\paragraph{(2)}

\begin{figure}[htb]
\begin{center}
\begin{tikzpicture}[scale=1.8, >=stealth]

  % --- 1. 座標定義 ---
  % 中央の面を triangle ABC (辺 AB を底辺) とし、
  % その周りに ACD, BCD, ABD を配置
  
  % 基準の三角形 ABC: AB = c, BC = a, CA = b (すべて異なる長さ)
  % ここでは描画例として a=sqrt(5), b=sqrt(6), c=sqrt(3) の比率を模式化
  \coordinate (A) at (0, 0);
  \coordinate (B) at (1.732, 0); % AB = c = sqrt(3)
  \coordinate (C) at (0.866, 2.291); % 頂点C (CA=b, CB=aから計算)

  % 各面を展開した頂点 D の座標 (D1, D2, D3)
  % 1. AC に関する合同三角形 (三角形 ACD_1 = 三角形 ABC)
  %    => AD1 = a, CD1 = c
  \coordinate (D1) at (-0.320, 2.766); % 頂点Cから、ACに関してBと反対側
  
  % 2. BC に関する合同三角形 (三角形 BCD_2 = 三角形 ABC)
  %    => BD2 = b, CD2 = c
  \coordinate (D2) at (2.052, 2.766); % 頂点Cから、BCに関してAと反対側
  
  % 3. AB に関する合同三角形 (三角形 ABD_3 = 三角形 ABC)
  %    => AD3 = b, BD3 = a
  \coordinate (D3) at (0.866, -2.291); % 辺ABに関してCと反対側

  % --- 2. 各面（合同な三角形）の塗りつぶし ---
  % 面の色を少し変えて区別
  \fill[blue!10] (A) -- (B) -- (C) -- cycle; % 中央 ABC
  \fill[red!10]  (A) -- (C) -- (D1) -- cycle; % ACD1
  \fill[green!10] (B) -- (C) -- (D2) -- cycle; % BCD2
  \fill[yellow!15] (A) -- (B) -- (D3) -- cycle; % ABD3

  % --- 3. 辺の描画と長さの表示 ---
  % 中央の面 (実線)
  \draw[thick] (A) -- (B) node[midway, below] {$c_1$};
  \draw[thick] (B) -- (C) node[midway, below right, yshift=1pt] {$a_2$};
  \draw[thick] (C) -- (A) node[midway, below left, yshift=1pt] {$b_1$};

  % 周囲の面 (折り目となる中央の辺は実線、外枠は太めの実線で区別)
  % 面 ACD1
  \draw[ultra thick] (A) -- (D1) node[midway, above left] {$a_1$};
  \draw[ultra thick] (C) -- (D1) node[midway, above] {$c_2$};
  
  % 面 BCD2
  \draw[ultra thick] (B) -- (D2) node[midway, above right] {$b_2$};
  \draw[ultra thick] (C) -- (D2) node[midway, above] {$c_2$};
  
  % 面 ABD3
  \draw[ultra thick] (A) -- (D3) node[midway, below left] {$a_1$};
  \draw[ultra thick] (B) -- (D3) node[midway, below right] {$b_2$};

  % --- 4. 頂点名とマーク ---
  % 中央の頂点
  \filldraw (A) circle(1.0pt) node[below left] {$A$};
  \filldraw (B) circle(1.0pt) node[below right] {$B$};
  \filldraw (C) circle(1.0pt) node[above, yshift=2pt] {$C$};

  % 展開された頂点 D
  \filldraw[black!60] (D1) circle(0.9pt) node[above left, black] {$D$};
  \filldraw[black!60] (D2) circle(0.9pt) node[above right, black] {$D$};
  \filldraw[black!60] (D3) circle(0.9pt) node[below, black] {$D$};

  % --- 5. 対辺のペアであることを示すマーク (同じ長さ c の辺) ---
  % AB=CD を示す二重線マーク
  \draw[thick, double] ($(A)!0.5!(B)$)+(-0.06,-0.12) -- +(-0.06,0.12);
  \draw[thick, double] ($(C)!0.5!(D1)$)+(-0.06,-0.10) -- +(-0.06,0.14);
  \draw[thick, double] ($(C)!0.5!(D2)$)+( 0.06,-0.10) -- +( 0.06,0.14);

\end{tikzpicture}
\end{center}
\caption{四面体の展開図と各辺の長さ}
\end{figure}

6つの辺の長さを$AB=c_1,\,CD=c_2,\,AC=b_1,\,BD=b_2,\,AD=a_1,\,BC=a_2$とおく．4つの面はそれぞれ3辺の組
\begin{align}
\triangle ABC:(a_2,b_1,c_1),\quad\triangle ABD:(a_1,b_2,c_1),\quad\triangle ACD:(a_1,b_1,c_2),\quad\triangle BCD:(a_2,b_2,c_2)
\end{align}
を持つ．全ての面が互いに合同なので，これらの3辺の組はすべて等しい：
\begin{align}
\{a_2,b_1,c_1\}=\{a_1,b_2,c_1\}=\{a_1,b_1,c_2\}=\{a_2,b_2,c_2\}.
\end{align}
$\triangle ABC$と$\triangle ABD$は辺$AB$（長さ$c_1$）を共有しているから残りの2辺が一致する：
\begin{align}
\{a_2,b_1\}=\{a_1,b_2\}.
\end{align}
同様に，$\triangle ABC,\triangle ACD$は辺$AC$を共有するので$\{a_2,c_1\}=\{a_1,c_2\}$，$\triangle ABD,\triangle ACD$は辺$AD$を共有するので$\{b_2,c_1\}=\{b_1,c_2\}$．これら3つの関係を同時にみたす組合せは
\begin{align}
a_1=a_2,\qquad b_1=b_2,\qquad c_1=c_2
\end{align}
のみであるから，
\begin{align}
|\overrightarrow{AC}|=|\overrightarrow{BD}|,\qquad |\overrightarrow{BC}|=|\overrightarrow{AD}|,\qquad|\overrightarrow{AB}|=|\overrightarrow{CD}|
\end{align}
である．よって題意は示された．

\paragraph{(3)}

(2)より対辺が等しい四面体なので，$AB=CD=\sqrt{3}$，$BC=AD=\sqrt{5}$，$CA=BD=\sqrt{6}$である．このような（対辺が等しい）四面体は直方体の4頂点として実現できる．直方体の辺の長さを$x,y,z$とし
\begin{align}
A=(0,0,0),\ B=(x,y,0),\ C=(x,0,z),\ D=(0,y,z)
\end{align}
とおくと
\begin{align}
AB=CD=\sqrt{x^2+y^2},\quad CA=BD=\sqrt{x^2+z^2},\quad BC=AD=\sqrt{y^2+z^2}
\end{align}
となり，与えられた長さから
\begin{align}
x^2+y^2=3,\qquad y^2+z^2=5,\qquad x^2+z^2=6.
\end{align}
とすれば良い．実際辺々を加えると$2(x^2+y^2+z^2)=14$，すなわち$x^2+y^2+z^2=7$だから
\begin{align}
 z^2=7-3=4,\qquad x^2=7-5=2,\qquad y^2=7-6=1 \\
\therefore
 x=\sqrt{2},\ y=1,\ z=2 \label{eq:2}
\end{align}
とすれば実現できる．

\begin{figure}[htb]
\begin{center}
\begin{tikzpicture}[
    scale=2.0,
    >=stealth,
    % 3D 投影ベクトルの定義 (x: 斜め手前, y: 右, z: 上)
    x={(-0.8cm, -0.6cm)},
    y={(0.8cm, -0.4cm)},
    z={(0cm, 1.0cm)}
    %x={(-0.65cm,-0.35cm)},
    %y={(0.95cm,-0.10cm)},
    %z={(0cm,1.15cm)}
  ]

  % --- 1. 3辺の長さ (x = sqrt(2) ≒ 1.414, y = 1.0, z = 2.0) ---
  \pgfmathsetmacro{\lx}{1.414}
  \pgfmathsetmacro{\ly}{1.0}
  \pgfmathsetmacro{\lz}{2.0}

  % --- 2. 直方体の8頂点の定義 ---
  \coordinate (V000) at (0, 0, 0);       % A
  \coordinate (Vx00) at (\lx, 0, 0);
  \coordinate (Vxy0) at (\lx, \ly, 0);   % B
  \coordinate (V0y0) at (0, \ly, 0);
  \coordinate (V00z) at (0, 0, \lz);
  \coordinate (Vx0z) at (\lx, 0, \lz);   % C
  \coordinate (Vxyz) at (\lx, \ly, \lz);
  \coordinate (V0yz) at (0, \ly, \lz);   % D

  % 四面体 ABCD の頂点
  \coordinate (A) at (V000);
  \coordinate (B) at (Vxy0);
  \coordinate (C) at (Vx0z);
  \coordinate (D) at (V0yz);

  % 中点 K(AB), L(BC), N(DA), P(CA)
  \coordinate (K) at ($(A)!0.5!(B)$);
  \coordinate (L) at ($(B)!0.5!(C)$);
  \coordinate (N) at ($(D)!0.5!(A)$);
  \coordinate (P) at ($(C)!0.5!(A)$);

  % --- 3. 直方体の枠線 (薄いグレー) ---
  % 奥側の見えない線（破線）
  \draw[dashed, gray!40, thin] (V000) -- (Vx00);
  \draw[dashed, gray!40, thin] (V000) -- (V0y0);
  \draw[dashed, gray!40, thin] (V000) -- (V00z);

  % 手前・外枠の線（実線）
  \draw[gray!50, thin] (Vx00) -- (Vxy0) -- (V0y0);
  \draw[gray!50, thin] (V00z) -- (Vx0z) -- (Vxyz) -- (V0yz) -- cycle;
  \draw[gray!50, thin] (Vx00) -- (Vx0z);
  \draw[gray!50, thin] (Vxy0) -- (Vxyz);
  \draw[gray!50, thin] (V0y0) -- (V0yz);

  % 辺の長さラベル
  \node[below left, font=\scriptsize, gray!80!black] at ($(Vx00)!0.5!(Vxy0)$) {$x=\sqrt{2}$};
  \node[below right, font=\scriptsize, gray!80!black] at ($(Vxy0)!0.5!(V0y0)$) {$y=1$};
  \node[left, font=\scriptsize, gray!80!black] at ($(V00z)!0.5!(V0yz)$) {$z=2$};

  % --- 4. 四面体 PKLN の描画（着色ハイライト） ---
  \fill[red!25, opacity=0.7] (P) -- (K) -- (L) -- cycle;
  \fill[red!35, opacity=0.7] (P) -- (L) -- (N) -- cycle;
  \fill[red!20, opacity=0.7] (P) -- (K) -- (N) -- cycle;
  
  \draw[thick, red!80!black] (P) -- (K) -- (L) -- cycle;
  \draw[thick, red!80!black] (P) -- (N) -- (L);
  \draw[thick, red!80!black, dashed] (K) -- (N);

  % --- 5. 四面体 ABCD の辺の描画 ---
  % 奥の線 (破線)
  \draw[thick, blue!80!black, dashed] (A) -- (B);
  \draw[thick, blue!80!black, dashed] (A) -- (C);
  \draw[thick, blue!80!black, dashed] (A) -- (D);

  % 手前の線 (太い実線)
  \draw[ultra thick, blue!90!black] (B) -- (C) -- (D) -- cycle;

  % --- 6. 各頂点のプロットとラベル ---
  % 四面体 ABCD
  \filldraw[blue!90!black] (A) circle (1.2pt) node[below right, font=\small] {$A$};
  \filldraw[blue!90!black] (B) circle (1.2pt) node[below right, font=\small] {$B$};
  \filldraw[blue!90!black] (C) circle (1.2pt) node[above left, font=\small] {$C$};
  \filldraw[blue!90!black] (D) circle (1.2pt) node[above right, font=\small] {$D$};

  % 中点 PKLN
  \filldraw[red!80!black] (K) circle (1.0pt) node[below, font=\scriptsize] {$K$};
  \filldraw[red!80!black] (L) circle (1.0pt) node[right, font=\scriptsize] {$L$};
  \filldraw[red!80!black] (N) circle (1.0pt) node[left, font=\scriptsize] {$N$};
  \filldraw[red!80!black] (P) circle (1.0pt) node[above left, font=\scriptsize] {$P$};

\end{tikzpicture}
\end{center}
\caption{直方体に埋め込まれた四面体 $ABCD$（青）と四面体 $PKLN$（赤）}
\end{figure}

四面体$PKLN$の体積は対称性より四面体$ABCD$の面積の$\dfrac{1}{8}$である．四面体$ABCD$の体積は$x,y,z$を$3$辺とする直方体から体積の同じ4つの三角錐を取り除いたもので，その体積は
\begin{align}
 \text{四面体}ABCD = xyz - 4 \cdot \dfrac{xyz}{6} = \dfrac{xyz}{3}
\end{align}
であるから，四面体$PKLN$の体積は
\begin{align}
 \text{四面体}PKLN = \dfrac{\text{四面体}ABCD}{8} = \dfrac{xyz}{24} = \dfrac{\sqrt{2}}{12}
\end{align}
である．

{\bf [別解1]}
% TODO :: 元の手書き回答の解法を追加

{\bf [別解2]}
行列式の知識があれば以下のように計算するのが一番機械的だろう．

\cref{eq:2}以下，$K,L,N,P$の座標はそれぞれ$AB,BC,DA,AC$の中点だから
\begin{align}
K=\Bigl(\dfrac{\sqrt{2}}{2},\dfrac{1}{2},0\Bigr),\quad L=\Bigl(\sqrt{2},\dfrac{1}{2},1\Bigr),\quad N=\Bigl(0,\dfrac{1}{2},1\Bigr),\quad P=\Bigl(\dfrac{\sqrt{2}}{2},0,1\Bigr) \label{eq:3}
\end{align}
となる．
$P$を基準にすると
\begin{align}
\overrightarrow{PK}=\Bigl(0,\dfrac{1}{2},-1\Bigr),\quad \overrightarrow{PL}=\Bigl(\dfrac{\sqrt{2}}{2},\dfrac{1}{2},0\Bigr),\quad\overrightarrow{PN}=\Bigl(-\dfrac{\sqrt{2}}{2},\dfrac{1}{2},0\Bigr)
\end{align}
であるから，四面体$PKLN$の体積は
\begin{align}
V=\dfrac{1}{6}\bigl|\det(\overrightarrow{PK},\overrightarrow{PL},\overrightarrow{PN})\bigr|.
\end{align}
ここで
\begin{align}
\det(\overrightarrow{PK},\overrightarrow{PL},\overrightarrow{PN})
=\begin{vmatrix}0&\dfrac{1}{2}&-1\\[2pt]\dfrac{\sqrt{2}}{2}&\dfrac{1}{2}&0\\[2pt]-\dfrac{\sqrt{2}}{2}&\dfrac{1}{2}&0\end{vmatrix}
=(-1)\times\Bigl(\dfrac{\sqrt{2}}{2}\cdot\dfrac{1}{2}-\dfrac{1}{2}\cdot\Bigl(-\dfrac{\sqrt{2}}{2}\Bigr)\Bigr)=(-1)\times\dfrac{\sqrt{2}}{2}=-\dfrac{\sqrt{2}}{2}
\end{align}
（第3列で展開した）．よって
\begin{align}
V=\dfrac{1}{6}\cdot\dfrac{\sqrt{2}}{2}=\dfrac{\sqrt{2}}{12}.
\end{align}


\end{document}
